\documentclass[12pt, twoside]{article}
\input{../preamble}

\fancyhead[LO]{La Scuola d'Italia / Dr.\ Huson / IB Math \\* 26 March 2026}
\fancyhead[RO]{Markscheme}
\fancyfoot[C]{\thepage}

\begin{document}
\subsubsection*{5.7 Classwork: Calculus Practice --- Markscheme}

\begin{enumerate}[itemsep=0.2cm]

%--- Problem 1: f(x) = x³ - 4x + 1 [9 marks] ---
\item Let $f(x) = x^3 - 4x + 1$.

\begin{enumerate}[itemsep=0.3cm]
  \item Find $f'(x)$. \hfill [2]

  \begin{minipage}[t]{0.62\textwidth}
  $f'(x) = 3x^2 - 4$
  \end{minipage}%
  \begin{minipage}[t]{0.35\textwidth}\raggedleft
  \textbf{(A1)(A1)}
  \end{minipage}

  \medskip
  \item Find the gradient of the graph of $f$ at $x = 2$. \hfill [2]

  \begin{minipage}[t]{0.62\textwidth}
  $f'(2) = 3(4) - 4$\\[0.15cm]
  $= 8$
  \end{minipage}%
  \begin{minipage}[t]{0.35\textwidth}\raggedleft
  \textbf{(M1)}\\[0.15cm]
  \textbf{(A1)}
  \end{minipage}

  \medskip
  \item Find the equation of the tangent line at $x = 2$. \hfill [3]

  \begin{minipage}[t]{0.62\textwidth}
  $f(2) = 8 - 8 + 1 = 1$\\[0.15cm]
  $y - 1 = 8(x - 2)$\\[0.15cm]
  $y = 8x - 15$
  \end{minipage}%
  \begin{minipage}[t]{0.35\textwidth}\raggedleft
  \textbf{(A1)}\\[0.15cm]
  \textbf{(M1)}\\[0.15cm]
  \textbf{(A1)}
  \end{minipage}

  \medskip
  \item Find the coordinates of the points where the gradient of $f$ is zero. \hfill [2]

  \begin{minipage}[t]{0.62\textwidth}
  $3x^2 - 4 = 0 \Rightarrow x = \pm\dfrac{2}{\sqrt{3}}
    = \pm\dfrac{2\sqrt{3}}{3}$\\[0.3cm]
  $\left(-\dfrac{2\sqrt{3}}{3},\; 1 + \dfrac{16\sqrt{3}}{9}\right)$
  and
  $\left(\dfrac{2\sqrt{3}}{3},\; 1 - \dfrac{16\sqrt{3}}{9}\right)$
  \end{minipage}%
  \begin{minipage}[t]{0.35\textwidth}\raggedleft
  \textbf{(A1)}\\[0.8cm]
  \textbf{(A1)}
  \end{minipage}
\end{enumerate}

\bigskip

%--- Problem 2: f(x) = 6x² - 2x [10 marks] ---
\item Let $f(x) = 6x^2 - 2x$.

\begin{enumerate}[itemsep=0.3cm]
  \item Find $\displaystyle\int (6x^2 - 2x)\,dx$. \hfill [3]

  \begin{minipage}[t]{0.62\textwidth}
  $\displaystyle\int (6x^2 - 2x)\,dx = 2x^3 - x^2 + c$
  \end{minipage}%
  \begin{minipage}[t]{0.35\textwidth}\raggedleft
  \textbf{(A1)(A1)(A1)}
  \end{minipage}

  \medskip
  \item Find the area of the shaded region. \hfill [3]

  \begin{minipage}[t]{0.62\textwidth}
  $\displaystyle\int_1^2 (6x^2 - 2x)\,dx
    = \left[2x^3 - x^2\right]_1^2$\\[0.2cm]
  $= (16 - 4) - (2 - 1)$\\[0.15cm]
  $= 12 - 1 = 11$
  \end{minipage}%
  \begin{minipage}[t]{0.35\textwidth}\raggedleft
  \textbf{(M1)}\\[0.2cm]
  \textbf{(A1)}\\[0.15cm]
  \textbf{(A1)}
  \end{minipage}

  \medskip
  \item Find $g(x)$ given $g'(x) = 6x^2 - 2x$ and $g(1) = 7$. \hfill [4]

  \begin{minipage}[t]{0.62\textwidth}
  $g(x) = 2x^3 - x^2 + c$\\[0.15cm]
  $g(1) = 2 - 1 + c = 7$\\[0.15cm]
  $c = 6$\\[0.15cm]
  $g(x) = 2x^3 - x^2 + 6$
  \end{minipage}%
  \begin{minipage}[t]{0.35\textwidth}\raggedleft
  \textbf{(M1)}\\[0.15cm]
  \textbf{(M1)}\\[0.15cm]
  \textbf{(A1)}\\[0.15cm]
  \textbf{(A1)}
  \end{minipage}
\end{enumerate}

\newpage

%--- Problem 3: f(x) = (2x + 1)eˣ [12 marks] ---
\item Let $f(x) = (2x + 1)e^x$.

\begin{enumerate}[itemsep=0.3cm]
  \item Show that $f'(x) = (2x + 3)e^x$. \hfill [3]

  \begin{minipage}[t]{0.62\textwidth}
  Product rule: $f'(x) = 2 \cdot e^x + (2x + 1) \cdot e^x$\\[0.15cm]
  $= e^x\big(2 + 2x + 1\big)$\\[0.15cm]
  $= (2x + 3)e^x$
  \end{minipage}%
  \begin{minipage}[t]{0.35\textwidth}\raggedleft
  \textbf{(M1)}\\[0.15cm]
  \textbf{(A1)}\\[0.15cm]
  \textbf{(A1)} \textbf{(AG)}
  \end{minipage}

  \medskip
  \item Find the equation of the tangent at $x = 0$. \hfill [3]

  \begin{minipage}[t]{0.62\textwidth}
  $f(0) = (1)(1) = 1$\\[0.15cm]
  $f'(0) = (3)(1) = 3$\\[0.15cm]
  $y = 3x + 1$
  \end{minipage}%
  \begin{minipage}[t]{0.35\textwidth}\raggedleft
  \textbf{(A1)}\\[0.15cm]
  \textbf{(A1)}\\[0.15cm]
  \textbf{(A1)}
  \end{minipage}

  \medskip
  \item Find the coordinates of $A$, where the tangent meets the $x$-axis. \hfill [2]

  \begin{minipage}[t]{0.62\textwidth}
  $0 = 3x + 1 \Rightarrow x = -\dfrac{1}{3}$\\[0.2cm]
  $A = \left(-\dfrac{1}{3},\; 0\right)$
  \end{minipage}%
  \begin{minipage}[t]{0.35\textwidth}\raggedleft
  \textbf{(M1)}\\[0.2cm]
  \textbf{(A1)}
  \end{minipage}

  \medskip
  \item Find the $x$-coordinate of the local minimum. Justify. \hfill [4]

  \begin{minipage}[t]{0.62\textwidth}
  $f'(x) = 0 \Rightarrow (2x + 3)e^x = 0$\\[0.15cm]
  $e^x > 0$ for all $x$, so $2x + 3 = 0$\\[0.15cm]
  $x = -\dfrac{3}{2}$\\[0.2cm]
  $f''(x) = (2x + 5)e^x$\\[0.15cm]
  $f''\!\left(-\dfrac{3}{2}\right) = 2 \cdot e^{-3/2} > 0$\\[0.15cm]
  Therefore it is a minimum.
  \end{minipage}%
  \begin{minipage}[t]{0.35\textwidth}\raggedleft
  \textbf{(M1)}\\[0.15cm]
  ~\\[0.15cm]
  \textbf{(A1)}\\[0.2cm]
  \textbf{(M1)}\\[0.15cm]
  ~\\[0.15cm]
  \textbf{(R1)}
  \end{minipage}
\end{enumerate}

\bigskip

%--- Problem 4: Kinematics [14 marks] ---
\item $v(t) = t^2 - 5t + 4$, \quad $s(0) = 2$.

\begin{enumerate}[itemsep=0.3cm]
  \item Write down the velocity when $t = 0$. \hfill [1]

  \begin{minipage}[t]{0.62\textwidth}
  $v(0) = 0 - 0 + 4 = 4$\;m\,s$^{-1}$
  \end{minipage}%
  \begin{minipage}[t]{0.35\textwidth}\raggedleft
  \textbf{(A1)}
  \end{minipage}

  \medskip
  \item Find the times when the particle is at rest. \hfill [3]

  \begin{minipage}[t]{0.62\textwidth}
  $t^2 - 5t + 4 = 0$\\[0.15cm]
  $(t - 1)(t - 4) = 0$\\[0.15cm]
  $t = 1$ and $t = 4$
  \end{minipage}%
  \begin{minipage}[t]{0.35\textwidth}\raggedleft
  \textbf{(M1)}\\[0.15cm]
  \textbf{(A1)}\\[0.15cm]
  \textbf{(A1)}
  \end{minipage}

  \medskip
  \item Find the acceleration when $t = 3$. \hfill [3]

  \begin{minipage}[t]{0.62\textwidth}
  $a(t) = v'(t) = 2t - 5$\\[0.15cm]
  $a(3) = 6 - 5$\\[0.15cm]
  $= 1$\;m\,s$^{-2}$
  \end{minipage}%
  \begin{minipage}[t]{0.35\textwidth}\raggedleft
  \textbf{(M1)(A1)}\\[0.15cm]
  ~\\[0.15cm]
  \textbf{(A1)}
  \end{minipage}

  \medskip
  \item Find $s(t)$. \hfill [4]

  \begin{minipage}[t]{0.62\textwidth}
  $s(t) = \displaystyle\int v\,dt = \frac{t^3}{3} - \frac{5t^2}{2} + 4t + c$\\[0.3cm]
  $s(0) = 2 \Rightarrow c = 2$\\[0.15cm]
  $s(t) = \dfrac{t^3}{3} - \dfrac{5t^2}{2} + 4t + 2$
  \end{minipage}%
  \begin{minipage}[t]{0.35\textwidth}\raggedleft
  \textbf{(M1)(A1)}\\[0.3cm]
  \textbf{(M1)}\\[0.15cm]
  \textbf{(A1)}
  \end{minipage}

  \medskip
  \item Find the total distance in the first 4 seconds. \hfill [3]

  \begin{minipage}[t]{0.62\textwidth}
  $v$ changes sign at $t = 1$ and $t = 4$\\[0.15cm]
  $\displaystyle\int_0^1 v\,dt
    = \left[\frac{t^3}{3} - \frac{5t^2}{2} + 4t\right]_0^1
    = \frac{1}{3} - \frac{5}{2} + 4 = \frac{11}{6}$\\[0.3cm]
  $\displaystyle\int_1^4 v\,dt
    = \left[\frac{t^3}{3} - \frac{5t^2}{2} + 4t\right]_1^4
    = \frac{-8}{3} - \frac{11}{6} = \frac{-9}{2}$\\[0.3cm]
  Total distance $= \dfrac{11}{6} + \dfrac{9}{2}
    = \dfrac{11}{6} + \dfrac{27}{6} = \dfrac{38}{6} = \dfrac{19}{3}$\;m
  \end{minipage}%
  \begin{minipage}[t]{0.35\textwidth}\raggedleft
  \textbf{(M1)}\\[0.15cm]
  ~\\[0.3cm]
  \textbf{(A1)}\\[0.6cm]
  \textbf{(A1)}
  \end{minipage}
\end{enumerate}

\end{enumerate}
\end{document}
